\documentclass[10pt,a4paper]{article}

\usepackage{fontspec}
\usepackage{polyglossia}
\setdefaultlanguage{russian}
\setotherlanguage{english}
\setmainfont{PT Sans}
\setmonofont[Scale=MatchLowercase]{PT Mono}
\newfontfamily\cyrillicfonttt[Scale=MatchLowercase]{PT Mono}

\usepackage[a4paper,top=1.0cm,bottom=1.0cm,left=1.3cm,right=1.3cm]{geometry}
\usepackage{titlesec}
\usepackage{enumitem}
\usepackage{xcolor}
\usepackage[hidelinks]{hyperref}
\definecolor{rule}{HTML}{333333}
\titleformat{\section}{\normalsize\bfseries}{}{0em}{\MakeUppercase}[{\color{rule}\titlerule[0.6pt]}]
\titlespacing*{\section}{0pt}{6pt}{3pt}
\setlist[itemize]{leftmargin=1.15em,labelsep=0.5em,itemsep=1.5pt,topsep=2pt,parsep=0pt}
\renewcommand{\labelitemi}{\textbullet}
\setlength{\parindent}{0pt}
\setlength{\parskip}{0pt}
\pagestyle{empty}
\linespread{1.02}
\tolerance=1500
\emergencystretch=2.5em
\newcommand{\role}[2]{\textbf{#1}\hfill #2\par}

\begin{document}
\begin{center}
  {\LARGE\bfseries Михайловский Илья Евгеньевич}\\[2pt]
  {\large DevOps / Platform Engineer}\\[3pt]
  Москва \textperiodcentered{} офис, гибрид или удалённо \textperiodcentered{} английский C1\\[2pt]
  Telegram: \href{https://t.me/Iluuuaaa}{@Iluuuaaa} \textperiodcentered{}
  \href{mailto:ilyamihailovsy@gmail.com}{ilyamihailovsy@gmail.com} \textperiodcentered{}
  \href{https://github.com/iluuua}{github.com/iluuua}
\end{center}

\section{О себе}
Разрабатываю backend и сопровождаю сервисы на Linux VPS. Автоматизирую сборку, развёртывание
и проверку релизов, пишу инструменты на Python и Bash. Основал ReachFlow; работаю с
контейнерами, базами данных и сетевой диагностикой. Учусь в Московском политехе, выпуск в 2029 году.

\section{Навыки}
\begin{itemize}
  \item Linux и сети: Ubuntu, SSH, systemd, диагностика процессов и соединений;
        TCP/IP, DNS, TLS, HTTP/2, reverse proxy, CDN, Caddy, nginx
  \item Сборка и эксплуатация: Docker, Docker Compose, Git, GitHub Actions,
        Bash, health checks, структурированные логи, резервное копирование и откат
  \item Разработка и данные: Python, FastAPI, asyncio, pytest, Ruff, PostgreSQL, SQL, Redis
  \item Kubernetes: базовое знакомство с манифестами Deployment, CronJob и Secret
\end{itemize}

\section{Опыт работы}
\role{ReachFlow, Founder}{январь 2026 - настоящее время}
\textit{B2B SaaS \textperiodcentered{} \href{https://reachflow.tech}{reachflow.tech}}
Основатель продукта; отвечаю за backend, инфраструктуру и выпуск изменений.
\begin{itemize}
  \item Настроил развёртывание API, шлюза, фоновых процессов и PostgreSQL через Docker Compose
        на Linux VPS. Использую Caddy/nginx для маршрутизации запросов и TLS.
  \item Настроил тесты и проверку Python-кода в GitHub Actions, деплой через webhooks и systemd.
        Устранил потерю событий при одновременных пушах: ожидание блокировки и таймер сверки версий.
  \item До миграций выполняю дамп PostgreSQL, проверку целостности и репетицию восстановления
        на отдельной базе. После релиза проверяю SHA образов и здоровье сервисов; предусмотрен откат.
  \item Разбирал зависание webhook receiver, ошибки доступа Docker в systemd и сбои reverse proxy.
        Использую журнал сервисов, проверки здоровья и структурированные логи.
\end{itemize}
\vspace{2pt}

\role{Delnovo, Chief Technology Officer}{июль 2023 - май 2026}
\begin{itemize}
  \item Разрабатывал ботов для поддержки и продаж: Python, PostgreSQL, интеграции с Битрикс24,
        запуск в Docker. По оценке заказчика бот-консультант автоматизировал около 80\%
        типовых консультаций.
\end{itemize}
\vspace{2pt}

\role{NEAR AI, Solutions Engineer}{июнь 2023 - июль 2024}
\begin{itemize}
  \item Решал алгоритмические задачи на C++ для обучения децентрализованного ИИ,
        оптимизировал время и память, проверял код участников международной команды.
\end{itemize}

\section{Проекты}
\role{ReachVPS, инфраструктура сервисов}{2026}
Работал с FastAPI, PostgreSQL, Redis, фоновым scheduler и Caddy в Docker Compose.
Проверял DNS, TLS, HTTP/2, целостность POST-запросов и потоковую передачу через Selectel CDN.

\vspace{3pt}
\role{Лидеры цифровой трансформации 2026, задача №5}{2026}
Капитан команды xixi-square. Разработал алгоритм обнаружения препятствий по данным 3D-лидара
и интеграцию с ROS 2 Humble. Решение запускается в Docker; отдельный процесс обрабатывает
последний кадр и пропускает устаревшие при перегрузке. На синтетике организаторов
найдены все 8 препятствий.

\section{Образование}
\textbf{Московский политехнический университет}, <<Прикладная математика и информатика>>,
выпуск 2029. \textbf{НИУ ВШЭ, ФКН}, вольнослушатель курсов той же программы.
\textbf{Т-Образование}, <<Т-Поколение>>, курс <<Алгоритмы>>, параллель B, 2025.
\textbf{Deep Learning School} ФПМИ МФТИ, 2025.

\end{document}
